% Generated draft: requires mathematical and source review.
\documentclass[11pt]{article}
\usepackage[margin=25mm]{geometry}
\usepackage{fontspec}
\setmainfont{Latin Modern Roman}
\usepackage{amsmath,amssymb,mathtools,unicode-math}
\setmathfont{Latin Modern Math}
\usepackage{enumitem,needspace,xurl,hyperref}
\setlength{\parindent}{0pt}
\setlength{\parskip}{5pt}
\setlength{\emergencystretch}{3em}
\newcommand{\N}{\mathbb{N}}
\newcommand{\Z}{\mathbb{Z}}
\newcommand{\Q}{\mathbb{Q}}
\newcommand{\R}{\mathbb{R}}
\newcommand{\C}{\mathbb{C}}
\newcommand{\F}{\mathbb{F}}
\newcommand{\A}{\mathbb{A}}
\newcommand{\PP}{\mathbb{P}}
\newcommand{\E}{\mathbb{E}}
\newcommand{\eps}{\varepsilon}
\newcommand{\dd}{\,\mathrm{d}}
\begin{document}
\section*{UnsolvedMath ID: 30004380}
\textbf{Source problem code:} catalogue entry
\subsection{Definitions and mathematical statement}
Let \(D_d\subset\R^d\) be a bounded domain, where \(d\geq 1\) is an integer, and let \(k\) denote the smoothness order appearing in the source. Under the usual nonnegative-integer interpretation of \(k\), \(C^k(D_d)\) denotes functions with continuous partial derivatives of orders at most \(k\) on \(D_d\). The supplied record does not specify the admissible range of \(k\), the scalar field, whether derivatives must extend continuously to the boundary, or any norm or bounded subset of this function space.

The symbol \(\operatorname{INT}\) denotes the intended integration functional on \(D_d\). Its measure, normalization, and integrability assumptions are not specified. Retain the source notation
\[
E_n(D_d):=e_n(C^k(D_d),\operatorname{INT}),\qquad n\geq 1,
\]
for its integration error quantity. The record does not define \(e_n\), including the permitted information, algorithm class, cost convention, or error criterion; consequently this notation cannot yet be interpreted as a uniquely defined numerical sequence.

\textbf{Source questions, conditional on specifying these conventions.} For fixed \(d\) and \(k\), determine whether
\[
L_{d,k}(D_d):=\limsup_{n\to\infty}E_n(D_d)\,n^{k/d}
\]
depends on the bounded domain \(D_d\). Equivalently, after specifying the admissible family of domains and using the same integration and error conventions throughout, is
\[
L_{d,k}(D_d)=L_{d,k}(D'_d)
\]
for every pair of admissible domains \(D_d,D'_d\subset\R^d\)? For each admissible \(D_d\), is the limsup an actual limit, that is, does
\[
\lim_{n\to\infty}E_n(D_d)\,n^{k/d}
\]
exist and equal \(L_{d,k}(D_d)\)? The source does not specify whether finiteness of this limit is required.

The original statement also asks more broadly how integration complexity depends on \(d\) and \(D_d\). Domain regularity, volume normalization, function-class bounds, and the computational model must be recovered before either the asymptotic questions or this broader request has a fully precise mathematical formulation.
\subsection{Short English statement}
For smooth multivariate integration, the source asks whether the asymptotic error constant after scaling by \(n^{k/d}\) depends on the integration domain, and whether the scaled errors converge. The function-class bounds, integration convention, admissible domains, and meaning of the error quantity remain unspecified, so this is a conditional statement requiring mathematical review.
\subsection{Sources}
\begin{itemize}
\item Local UnsolvedMath registry, permanent ID \(30004380\), source problem code \texttt{catalogue entry}, in the supplied complete source record. Project registry locations: \texttt{data/unsolvedmath/problems.json} and \texttt{catalog/problems.json}.
\item Supplied report citation: \emph{New Perspectives and Computational Challenges in High Dimensions} (2021), Oberwolfach Reports, DOI \texttt{10.4171/owr/2020/6}, pp.\ 377--416. Source URL and supplied literature-source link, labeled ``Oberwolfach report'': \url{https://doi.org/10.4171/owr/2020/6}. The collection attributes the problem to an Oberwolfach workshop report; no individual proposer is supplied. The proposed year and year tag are 2020, while the supplied citation gives 2021.
\item The supplied registry labels the problem \texttt{open} and the statement \texttt{corrected\_verified}. Its verification note says that the official report was checked and that the formula and its two questions were restored in place of extraction artifacts. This is supplied provenance, not an independent source check for this draft.
\item Literature assessment dated exactly \texttt{2026-08-22}: the supplied curation states that known bounds determine the principal dimension-dependent rate in several smooth classes, while questions about the exact asymptotic constant, domain dependence, and existence of the limit remain open. This is supplied dated curation, not independently verified. The background explicitly requests independent verification of time-sensitive status and cited claims.
\end{itemize}
\end{document}
